This thesis represents an exploration of a connection between the method of self-adjoint extension and the phenomenon of anomalous symmetry breaking referred to scale symmetry in quantum mechanics. At first glance, these two concepts may appear completely independent, but as we delve into the heart of this thesis, you will come to understand the link that binds them in our specific problem.

Through rigorous analysis, critical insights, and a thorough examination of relevant literature, this work strives to elucidate how self-adjoint extensions can shed light on the intricate interplay between symmetries and quantum systems. We aim to provide a bridge between the seemingly abstract world of mathematical operators and the tangible consequences of symmetry breaking in the realm of quantum physics.

By embarking on this intellectual journey, you will gain a deeper appreciation for the profound impact of self-adjoint extensions on the understanding of anomalous symmetry breaking, and how this connection shapes the very foundations of quantum mechanics.

As you embark on the chapters that follow, we invite you to explore this territory of theoretical physics and uncover the relationship between self-adjoint extensions and the anomalies that can emerge when quantum mechanics meets symmetries.

\vspace{4mm}
\noindent
Introducing the contents of the essay we can start saying that it is commonly taught that an attractive potential vanishing at infinity implies a set of bound states that are both discrete and infinite. 
However, we will demonstrate that this holds true only for certain potentials, as discussed in Section~\ref{section:generic_potentials}.

Our discussion will lead us to categorize attractive potentials that vanish at infinity into two groups, and this classification naturally leads to the definition of a special case, namely, the $-\alpha/x^2$ potential. 
This potential exhibits unique characteristics that set it apart from other potentials as we will show in Section~\ref{section:peculiarities}.

We will introduce and employ the method of self-adjoint extension, as described in Section~\ref{section:self_adjoint}, to address the issue that the naive formulation of the special case potential $-\alpha/x^2$ does not result in a Hamiltonian that is an Observable when considered in the context of Quantum Mechanics, as discussed in Section~\ref{section:xsqared_halfaxis}.
Additionally, we will provide a simpler introductory example of the mathematical method: the free particle on the half-axis in Section \ref{section:free_halfaxis}.

At the conclusion of this essay, we will observe that, upon elevating the Hamiltonian operator to the status of an Observable, the scale invariance of our problem has been unexpectedly broken. This anomaly will be examined from two perspectives: the boundary conditions and the bound states.

In conclusion, the essay will establish a connection between the problem discussed and the two-dimensional $\delta$-function potential problem commonly encountered in the literature on anomalies in Quantum Mechanics, as detailed in Section~\ref{section:delta_potential}.